\section{Kết luận và hướng phát triển}
\label{sec:conclusion}

Báo cáo đã \emph{nghiên cứu mô-đun phát hiện bất thường SyscallAD} của framework
DeSFAM trong ngữ cảnh phát hiện tấn công leo thang đặc quyền dựa trên lời gọi hệ
thống trong Kubernetes. Ứng với mục tiêu đề ra:

\begin{itemize}[leftmargin=1.8em]
  \item \textbf{Hiểu thiết kế.} Báo cáo đã hệ thống hóa đường ống của SyscallAD (cửa
        sổ trượt $\to$ véc-tơ 43 chiều $\to$ VAE\,+\,iForest $\to$ hợp nhất $\to$
        ngưỡng) và làm rõ \emph{cơ chế} phát hiện leo thang đặc quyền: các syscall
        đặc quyền hiếm được mã hóa bằng bit hiện diện nhị phân (\texttt{disc}), bền
        với pha loãng, nằm xa phân phối lành tính mà VAE đã học nên cho điểm bất
        thường cao.
  \item \textbf{Tái lập \& bằng chứng.} Trên DongTing đầy đủ, lõi phát hiện đạt AUC
        tập hợp $0{,}835$ (cửa sổ) / $0{,}787$ (chuỗi), AP $0{,}990$, precision
        $0{,}999$, ổn định qua hạt giống ($\pm 0{,}000$). Vì DeSFAM không công bố bảng
        số theo mô hình trên DongTing, báo cáo \emph{không} tuyên bố ``đạt/vượt số
        liệu công bố'', chỉ tham chiếu định tính (AUC cùng tầm mức tổng hợp
        $\approx 0{,}94$; F1/recall thấp hơn do ngưỡng \texttt{p995} + bảng chữ cái
        hẹp). Thí nghiệm cắt bỏ cho thấy ở mức cửa sổ năng lực phát hiện phần lớn quy
        về \emph{hiện diện} syscall đặc quyền (chỉ-\texttt{disc} $0{,}875 \ge$ tập hợp
        $0{,}835$), còn giá trị riêng của mô hình bộc lộ ở \emph{mức chuỗi} ($0{,}787$
        so với $0{,}678$)---một \emph{bằng chứng sơ bộ} có giới hạn.
  \item \textbf{Nền tảng tái lập.} Toàn bộ quy trình được đóng gói bằng Docker và
        ràng buộc bằng kiểm thử ngang hàng, cho phép các nghiên cứu sau chạy lại và
        mở rộng một cách độc lập.
\end{itemize}

Về bài học phương pháp (độc lập với con số): với phát hiện bất thường, (a) việc chọn
lọc đặc trưng theo miền (giới hạn về tập syscall an ninh) định hình mạnh kết quả, và
(b) \emph{đơn vị đánh giá} (cửa sổ so với chuỗi) cùng \emph{tiêu chí đặt trước} quan
trọng ngang chất lượng mô hình---bỏ qua chúng dễ dẫn tới kết luận sai lệch.

\textbf{Hướng phát triển.} (1) Chọn ngưỡng tối ưu F1 và báo cáo đầy đủ đường cong
precision--recall; (2) phép đối chứng bảng chữ cái \texttt{SENSOR\_ALPHABET} 0/1 để
định lượng đóng góp của bước chọn lọc đặc trưng---bộ công cụ đã sẵn trong
\texttt{Experiment/}; (3) bổ sung đặc trưng thời gian khi có nguồn dữ liệu kèm dấu
thời gian theo syscall (điều DongTing thiếu); (4) mở rộng nguồn dữ liệu và dùng vết
khai thác CVE thực; (5) điều tra bất đối xứng mã hóa của
\texttt{setgid}/\texttt{unlinkat} (Mục~\ref{subsec:confound}); (6) (xa hơn) kết nối
khâu thực thi của DeSFAM để chuyển từ \emph{phát hiện} sang \emph{ngăn chặn}.
